
\SGASection{3}{Characterization of a sub-torus by its underlying set}
\SGARefTarget{sga3:xv:section:3}

\SGASubsection{3.1}{Statement of the theorem}
\SGARefTarget{sga3:xv:subsection:3:1}

\noindent\textbf{Notation.}---If \(X\) is a prescheme,
\(\operatorname{ens}(X)\) denotes the underlying set of \(X\). If
\(X\) and \(S'\) are two \(S\)-preschemes,
\[
  X_{S'}=X\times_S S',
  \qquad u:X_{S'}\longrightarrow X
\]
is the canonical morphism, and if \(E\) is a subset of
\(\operatorname{ens}(X)\), we denote by \(E_{S'}\) (or
\(E\times_S S'\)) the subset \(u^{-1}(E)\) of
\(\operatorname{ens}(X_{S'})\).

\SGAStatement{\SGAInlineRefTarget{sga3:xv:theorem:3-1}Theorem 3.1}{Let \(S\) be a locally noetherian
prescheme, let \(G\) be an \(S\)-group prescheme of finite type, let
\(q>0\) be an integer invertible on \(S\), and let \(E\) be a subset
of \(\operatorname{ens}(G)\). Consider the following assertions about
\(E\):
\begin{enumerate}
\item[\normalfont\SGAInlineRefTarget{sga3:xv:theorem:3-1:item:i}(i)] The set \(E\) is the underlying set of a
sub-torus \(T\) of \(G\).
\item[\normalfont\SGAInlineRefTarget{sga3:xv:theorem:3-1:item:ii}(ii)]
  \begin{enumerate}
  \item[\normalfont a)] For every point \(s\) of \(S\), there is a
  sub-torus \(T_s\) of \(G_s\) such that
  \[
    E\cap\operatorname{ens}(G_s)=E_s
  \]
  is the underlying set of \(T_s\).
  \item[\normalfont b)] As \(n\) ranges over the powers of \(q\),
  there is a coherent family (cf.\ \SGARefLink{sga3:xv:corollary:2-2}{2.2}) \(M_n\) of subgroup schemes
  of \(G\), of multiplicative type, such that for every point \(s\)
  of \(S\),
  \[
    (M_n)_s={}_nT_s.
  \]
  \end{enumerate}
\item[\normalfont\SGAInlineRefTarget{sga3:xv:theorem:3-1:item:iii}(iii)]
  \begin{enumerate}
  \item[\normalfont a)] As in \textup{\SGARefLink{sga3:xv:theorem:3-1:item:ii}{(ii)} a)}.
  \item[\normalfont b)] The set \(E\) is locally closed in
  \(\operatorname{ens}(G)\), and the dimension of the fibers of \(E\)
  over \(S\) is locally constant.
  \end{enumerate}
\item[\normalfont\SGAInlineRefTarget{sga3:xv:theorem:3-1:item:iv}(iv)]
  \begin{enumerate}
  \item[\normalfont a)] As in \textup{\SGARefLink{sga3:xv:theorem:3-1:item:ii}{(ii)} a)}.
  \item[\normalfont b)] For every \(S\)-scheme \(S'\) that is the
  spectrum of a complete discrete valuation ring with algebraically
  closed residue field, \(E_{S'}\) is the underlying set of a
  sub-torus of \(G_{S'}\).
  \end{enumerate}
\end{enumerate}
Then the following implications hold:
\begin{enumerate}
\item[\normalfont A)] \(\textup{(i)}\Longleftrightarrow
\textup{(ii)}\Longrightarrow\textup{(iii)}
\Longrightarrow\textup{(iv)}\).
\item[\normalfont B)] If \(G\) is separated over \(S\), then
\(\textup{(iii)}\Longleftrightarrow\textup{(iv)}\), and moreover
\(E\) is closed.
\item[\normalfont C)] Conditions \textup{\SGARefLink{sga3:xv:theorem:3-1:item:i}{(i)}}, \textup{\SGARefLink{sga3:xv:theorem:3-1:item:ii}{(ii)}}, and
\textup{\SGARefLink{sga3:xv:theorem:3-1:item:iii}{(iii)}} (and also \textup{\SGARefLink{sga3:xv:theorem:3-1:item:iv}{(iv)}} if \(G\) is separated over
\(S\)) are equivalent in the following two cases:
  \begin{enumerate}
  \item[\normalfont First case:]
    \begin{enumerate}
    \item[\normalfont a)] The prescheme \(S\) is reduced, or \(G\)
    is flat over \(S\); and
    \item[\normalfont b)] for every point \(s\) of \(S\), \(T_s\)
    is a central torus of \(G_s\).
    \end{enumerate}
  \item[\normalfont Second case:] \(S\) is normal.
  \end{enumerate}
\end{enumerate}
Moreover, in the two cases above, the torus \(T\) with underlying
set \(E\) is unique.}
\SGARefTarget{sga3:xv:theorem:3:1}

\SGAStatement{\SGAInlineRefTarget{sga3:xv:remark:3-2}Remarks 3.2}{\SGAInlineRefTarget{sga3:xv:remark:3-2:item:a}\textup{a)} When \(S\) is reduced, there
is no need in \textup{\SGARefLink{sga3:xv:theorem:3-1:item:ii}{(ii)}} to assume that the family \(M_n\) is
coherent, since this condition follows from the other hypotheses.
Indeed, if \(m\) divides \(n\), the subgroups \({}_mM_n\) and \(M_m\)
are étale over \(S\), hence reduced, and have the same underlying
space; they therefore coincide.

\textup{b)} If \(S\) is no longer assumed normal, it is in general no
longer true that \(\textup{(iii)}\Rightarrow\textup{(i)}\), even when
\(S\) is reduced and geometrically unibranch and \(G\) is a smooth
group scheme over \(S\). Indeed, consider the Borel subgroup of
\(\operatorname{SL}_{2,S}\) consisting of matrices of the form
\[
  \begin{pmatrix}t&u\\0&t^{-1}\end{pmatrix},
\]
where \(S\) is the affine curve over a field \(k\) with ring
\[
  k[x,y]/(y^2-x^3).
\]
Define \(E\) as follows. Above the cusp of \(S\), where \(x=y=0\),
take the diagonal torus \(u=0\). Above the complementary open subset
\(x\ne0\), take the torus obtained from the diagonal torus by
conjugation by
\[
  \begin{pmatrix}1&y/x\\0&1\end{pmatrix}.
\]
The resulting set \(E\) satisfies \textup{\SGARefLink{sga3:xv:theorem:3-1:item:iii}{(iii)} a)}; it is also
closed in \(G\), and the reduced subscheme having \(E\) as its
underlying set has equations
\[
  \begin{cases}
    xu+y(t-t^{-1})=0,\\
    u^2-x(t-t^{-1})^2=0.
  \end{cases}
\]
It is therefore not a sub-torus of \(G\), since its fiber above the
cusp is not reduced.}
\SGARefTarget{sga3:xv:remarks:3:2}

\noindent\textbf{Plan of the proof of \SGARefLink{sga3:xv:theorem:3-1}{3.1}.} In a first part we shall
establish the following ``easy'' implications:
\[
\begin{tikzcd}[column sep=huge,row sep=large]
 & \textup{(ii)} \arrow[dr,Rightarrow] & \\
\textup{(i)} \arrow[ur,Rightarrow] \arrow[dr,Rightarrow] && \textup{(iv)} \\
 & \textup{(iii)} \arrow[ur,Rightarrow] &
\end{tikzcd}
\]
and
\[
  \textup{(iv)}
  \Longrightarrow
  [\,\textup{(iii)}\text{ and }E\text{ is closed}\,]
  \quad\text{if \(G\) is separated over \(S\)}.
\]
The more delicate implications will be proved in three stages:
\begin{enumerate}
\item[\normalfont I)] Reduction of
\(\textup{(iii)}\Rightarrow\textup{(i)}\), under the hypotheses of
C), first case, to the case in which \(S\) is normal.
\item[\normalfont II)] \(\textup{(iii)}\Rightarrow\textup{(ii)}\)
when \(S\) is normal.
\item[\normalfont III)] \(\textup{(ii)}\Rightarrow\textup{(i)}\).
\end{enumerate}

\SGASubsection{3.2}{Proof of the ``easy'' assertions in \SGARefLink{sga3:xv:theorem:3-1}{3.1}}
\SGARefTarget{sga3:xv:subsection:3:2}

The implications \(\textup{(i)}\Rightarrow\textup{(ii)}\) and
\(\textup{(i)}\Rightarrow\textup{(iii)}\) are trivial.

\medskip
\noindent\(\textup{(iii)}\Rightarrow\textup{(iv)}\).
After replacing \(S\) by \(S'\), we may assume that \(S\) is the
spectrum of a discrete valuation ring. Let \(t\) be the generic point
of \(S\), and let \(s\) be its closed point.

Since \(E\) is locally closed, there is a reduced subprescheme of
\(G\) with underlying space \(E\); we also denote it by \(E\). Its
generic fiber \(E_t\) is therefore the sub-torus \(T_t\) of \(G_t\),
of rank \(r\), by \textup{\SGARefLink{sga3:xv:theorem:3-1:item:iii}{(iii)} a)}. Let \(E'\) be the schematic
closure of \(E_t\) in \(E\). The prescheme \(E'\) is irreducible, and
its closed fiber \(E'_s\) is nonempty (it contains the unit section
of \(G_s\)); hence \(E'_s\) has dimension \(r\) (EGA IV~14.3.10).
But \(E'_s\) is a closed subscheme of \(E_s\), while \(E_s\) has
dimension \(r\) by \textup{\SGARefLink{sga3:xv:theorem:3-1:item:iii}{(iii)} b)} and is irreducible because it
has the same underlying space as \(T_s\). Thus \(E'_s\) and \(E_s\)
have the same underlying space, and consequently \(E'=E\). This
proves that \(E\) is flat over \(S\).

Let \(E''\) now be the schematic closure of \(E_t\) in \(G\). Then
\(E''\) is a subgroup prescheme of \(G\), flat over \(S\)
(Exposé \SGARefLink{sga3:viii:lemma:7:1}{VIII~7.1}), which contains \(E'\), and therefore \(E\). Its
closed fiber \(E''_s\) is an algebraic subgroup of \(G_s\) of
dimension \(r\) (\emph{loc.\ cit.}). Since \(T_s\) is a closed
irreducible subset of \(G_s\) of dimension \(r\), \(E_s\) has the
same underlying set as the identity component \((E''_s)^0\) of
\(E''_s\). Let \((E'')^0\) denote the ``identity component'' of
\(E''\), that is, the open subgroup of \(E''\) complementary to the
union of the irreducible components of \(E''_s\) that do not contain
the origin. Then
\[
  E=\operatorname{ens}\bigl((E'')^0\bigr).
\]
Since \(E\) and \(E''\) are reduced, one even has
\(E=(E'')^0\). Finally, \(E\) is a subgroup prescheme of \(G\), flat
and of finite type over \(S\), with connected fibers, hence separated
(Exposé VI\textsubscript{B}~\SGARefLink{sga3:xv:theorem:5-2}{5.2}); its generic fiber is the torus
\(T_t\), and its reduced closed fiber is the torus \(T_s\). It follows
that \(E\) is a torus (Exposé \SGARefLink{sga3:x:theorem:8-8}{X~8.8}).

\medskip
\noindent\(\textup{(ii)}\Rightarrow\textup{(iv)}\).
Again we may assume that \(S\) is the spectrum of a discrete valuation
ring, and we retain the notation introduced above. The schematic
closure \(T\) of \(T_t\) in \(G\) is a subgroup prescheme of \(G\),
flat over \(S\), and it contains \(M_n\) for every \(n\) that is a
power of \(q\). Consequently, its closed fiber is a closed subscheme
of \(G_s\) containing \((M_n)_s\) for every \(n\), hence containing
\(T_s\). For dimension reasons, the ``identity component'' \(T^0\)
of \(T\) has \(E\) as its underlying set, and the preceding argument
shows that \(T^0\) is a sub-torus of \(G\).

\medskip
\noindent\(\textup{(iv)}\Rightarrow
[\,\textup{(iii)}\text{ and }E\text{ is closed}\,]\), if \(G\) is
separated over \(S\). Let us show that \(E\) is
closed.\footnote{Editor's note: The implication
\(\textup{(iv)}\Rightarrow\textup{(iii)}\) will appear in the course
of the proof.} We first prove the following lemma.

\SGAStatement{Lemma 3.3}{If the conditions stated in \SGARefLink{sga3:xv:theorem:3-1}{3.1}
\textup{\SGARefLink{sga3:xv:theorem:3-1:item:iv}{(iv)}} are satisfied, \(E\) is a locally constructible subset
of \(\operatorname{ens}(G)\).}
\SGARefTarget{sga3:xv:lemma:3:3}

By the usual method, we are reduced to the case in which \(S\) is
noetherian and integral, with generic point \(\eta\). After shrinking
\(S\), we may assume (Exposé VI\textsubscript{B}~\SGARefLink{par:sga3-VIB-10.10}{10.10}) that there is
a subgroup scheme \(H\) of \(G\), flat over \(S\) and with connected
fibers, whose generic fiber \(H_\eta\) is \(T_\eta\). To prove \SGARefLink{sga3:xv:lemma:3:3}{Lemma
3.3} it is enough to show that \(\operatorname{ens}(H)=E\).

Let \(s\) be a point of \(S\). By EGA II~7.1.9 there is an
\(S\)-scheme \(S'\), the spectrum of a discrete valuation ring which
we may assume complete with algebraically closed residue field, whose
generic point \(t'\) maps to \(\eta\) and whose closed point \(s'\)
maps to \(s\). By \textup{\SGARefLink{sga3:xv:theorem:3-1:item:iv}{(iv)} b)}, there is a sub-torus \(T_{S'}\)
of \(G_{S'}\) having \(E_{S'}\) as its underlying space. The two
subgroup preschemes \(T_{S'}\) and
\(H_{S'}=H\times_S S'\) of \(G_{S'}\) are flat over \(S'\), have
connected fibers, and have the same generic fiber \(T_{t'}\).
Therefore both coincide with the identity component of the schematic
closure of \(T_{t'}\) in \(G_{S'}\). Hence
\(\operatorname{ens}(H_s)=E_s\), and therefore
\(\operatorname{ens}(H)=E\), proving the lemma.

Now that \(E\) is known to be locally constructible, to see that it
is closed it is enough to prove that every specialization in \(G\)
of a point of \(E\) again belongs to \(E\). By the usual technique,
we are reduced to the case in which \(S\) is the spectrum of a
complete discrete valuation ring with algebraically closed residue
field. In that case the sub-torus of \(G\) with underlying space
\(E\), supplied by \textup{\SGARefLink{sga3:xv:theorem:3-1:item:iv}{(iv)} b)}, is closed in \(G\), because
\(G\) is separated (Exposé \SGARefLink{sga3:viii:proposition:7:12}{VIII~7.12}).

\SGASubsection{3.3}{Continuation of the proof of \SGARefLink{sga3:xv:theorem:3-1}{3.1}}
\SGARefTarget{sga3:xv:subsection:3:3}

\noindent\textbf{I) Reduction of
\(\textup{(iii)}\Rightarrow\textup{(i)}\), case C, first case, to
the case in which \(S\) is normal.}

\smallskip
\noindent\textit{a) Reduction to the case in which \(S\) is affine
and reduced.}\quad
We therefore assume that, for every point \(s\) of \(S\), \(E_s\)
is the underlying space of a central sub-torus of \(G_s\). The
uniqueness of \(T\) then follows from Exposé \SGARefLink{sga3:ix:theorem:5:1:bis}{IX~5.1 bis}, and moreover
\(T\) will necessarily be central (\emph{loc.\ cit.}). In view of
this uniqueness, in order to prove the existence of \(T\) we may
suppose that \(S\) is noetherian and affine, with ring \(A\). If
\(S\) is not reduced, \(G\) is flat over \(S\) by hypothesis. By
\SGARefLink{sga3:xv:corollary:2:3}{2.3} it is then enough to solve the problem for \(S_{\mathrm{red}}\)
and \(G\times_S S_{\mathrm{red}}\). We may therefore suppose in
addition that \(S\) is reduced.

\smallskip
\noindent\textit{b) Reduction to the case in which \(A\) is an
algebra of finite type over \(\mathbb Z\).}\quad
We shall use the following lemma.

\SGAStatement{Lemma 3.4}{Let \(k\) be a field, \(G\) an algebraic
\(k\)-group, and \(E\) a subset of \(G\). Let \(k'/k\) be a field
extension, and suppose that \(E_{k'}\) is the underlying set of a
sub-torus \(T_{k'}\) of \(G_{k'}\). If \(T_{k'}\) is central in
\(G_{k'}\), or if \(k\) is perfect, then \(E\) is the underlying set
of a sub-torus \(T\) of \(G\).}
\SGARefTarget{sga3:xv:lemma:3:4}

\SGAProof{Proof}
By fpqc descent it is enough to prove that the two inverse images of
\(T_{k'}\) in \(G_{k''}\), where
\[
  k''=k'\otimes_k k',
\]
coincide. They have the same underlying set. If \(T_{k'}\) is
central, they therefore coincide by Exposé \SGARefLink{sga3:ix:theorem:5:1:bis}{IX~5.1 bis}. If \(k\) is
perfect, \(k''\) is reduced; the two inverse images, being smooth
over \(k''\), are reduced and consequently coincide.

\SGAStatement{\SGAInlineRefTarget{sga3:xv:remark:3-5}Remarks 3.5}{\textup{(1)} \SGARefLink{sga3:xv:lemma:3:4}{Lemma 3.4} applies in
particular when the residue fields of \(S\) are perfect.

\textup{\SGARefLink{sga3:xv:theorem:5:2}{(2)}} Under the hypotheses of the first case in \SGARefLink{sga3:xv:theorem:3-1}{Theorem 3.1}
C, for every \(S\)-scheme \(S'\) as in \SGARefLink{sga3:xv:theorem:3-1}{3.1} \textup{\SGARefLink{sga3:xv:theorem:3-1:item:iv}{(iv)} b)}, the
sub-torus of \(G_{S'}\) having \(E_{S'}\) as its underlying space is
central.}
\SGARefTarget{sga3:xv:remarks:3:5}

\SGAStatement{Lemma 3.6}{Let \(S\) be the projective limit of a
filtered inverse system of affine schemes \(S_i\), and let \(H\) be
an \(S\)-group of multiplicative type and of finite type. There
exist an index \(i\), a group \(H_i\) of multiplicative type over
\(S_i\), and an \(S\)-isomorphism
\[
  H_i\times_{S_i}S\xrightarrow{\sim}H.
\]
If \(H\) is isotrivial, \(H_i\) may be chosen isotrivial.}
\SGARefTarget{sga3:xv:lemma:3:6}

\SGAProof{Proof}
Since \(H\) is of finite presentation over \(S\), there are an index
\(i\), an \(S_i\)-group prescheme \(H_i\), of finite presentation
over \(S_i\), and an \(S\)-isomorphism
\[
  H_i\times_{S_i}S\xrightarrow{\sim}H
\]
(Exposé VI\textsubscript{B}~\SGARefLink{prop:sga3-VIB-10.16}{10.16}). After increasing \(i\), we may
suppose that \(H_i\) is of multiplicative type
(Exposé \SGARefLink{sga3:ix:proposition:2:11}{IX~2.11}). It remains to prove that \(H_i\) can be chosen
quasi-isotrivial, and isotrivial when \(H\) is isotrivial.

We may thus suppose that \(H\) is trivialized by an étale,
surjective morphism \(S'\to S\) of finite presentation. For \(i\)
sufficiently large, there are an \(S_i\)-prescheme \(S'_i\), an
étale, surjective morphism \(S'_i\to S_i\) of finite presentation,
and an \(S\)-isomorphism
\[
  S'_i\times_{S_i}S\xrightarrow{\sim}S'
\]
(EGA IV~17.16). For \(j\) sufficiently large set
\[
  S'_j=S'_i\times_{S_i}S_j,\qquad
  H'_j=H_j\times_{S_j}S'_j,\qquad
  H'=H\times_S S'.
\]
By the choice of \(S'\), there are a finitely generated abelian
group \(M\) and an \(S'\)-isomorphism of group schemes
\[
  D_{S'}(M)\xrightarrow{\sim}H'.
\]
Since the \(S'_i\) are quasi-compact and
\(S'=\varprojlim S'_i\), Exposé VI\textsubscript{B}~10 shows that
there are an index \(j\) and an \(S'_j\)-isomorphism of group
schemes
\[
  D_{S'_j}(M)\xrightarrow{\sim}H'_j.
\]
This says precisely that \(H_j\) is a quasi-isotrivial group of
multiplicative type. When \(H\) is isotrivial, one argues similarly
using a trivialization of \(H\) by a finite étale morphism
\(S'\to S\).

We can now carry out reduction b). The ring \(A\) of \(S\) is the
filtered direct limit of its subrings \(A_i\) of finite type over
\(\mathbb Z\). Put \(S_i=\operatorname{Spec}A_i\), and let
\[
  u_j:S\longrightarrow S_j,\qquad
  u_{jk}:S_k\longrightarrow S_j\quad(k\geq j)
\]
be the transition morphisms. Since \(S\) is noetherian and \(G\) is
of finite presentation over \(S\), there are an index \(i\) and an
\(S_i\)-group prescheme \(G_i\), of finite type over \(S_i\), such
that \(G\) is \(S\)-isomorphic to \(G_i\times_{S_i}S\). Likewise,
since \(E\) is locally closed in \(G\), we may suppose that there is
a locally closed subset \(E_i\) of \(G_i\) such that
\[
  E=E_i\times_{S_i}S
\]
(EGA IV~8.3.11). For \(j>i\), put
\[
  G_j=G_i\times_{S_i}S_j,\qquad
  E_j=E_i\times_{S_i}S_j,
\]
and let \(Q_j\) be the set of points \(s\) of \(S_j\) for which
\((E_j)_s\) is the underlying set of a central sub-torus of
\((G_j)_s\).

It follows from \SGARefLink{sga3:xv:lemma:3:4}{3.4} that
\[
  Q_k=u_{jk}^{-1}(Q_j)\quad(k\geq j),
\]
and by hypothesis
\[
  \operatorname{ens}(S)=u_j^{-1}(Q_j)\quad(j\geq i).
\]
I further claim that \(Q_j\) is ind-constructible
(EGA IV~1.9.4). Indeed, since \(S_j\) is noetherian, it is enough
(EGA IV~1.9.10) to see the following. Let \(S\) be a noetherian
integral scheme with generic point \(\eta\), and suppose that
\(E_\eta\) is the underlying set of a central sub-torus \(T_\eta\)
of \(G_\eta\). Then there is a neighborhood \(U\) of \(\eta\) such
that \(E_s\) has the same property for every \(s\in U\). After
shrinking \(S\), we may suppose that the center \(Z\) of \(G\) is
representable and that there is a subgroup scheme \(T\) of \(G\)
whose generic fiber is \(T_\eta\)
(Exposé VI\textsubscript{B}, \S10). The two locally closed subsets
\(\operatorname{ens}(T)\) and \(E\) of \(\operatorname{ens}(G)\)
coincide on the generic fiber. We can therefore find a neighborhood
\(U\) of \(\eta\) such that, above \(U\),
\[
  \operatorname{ens}(T)=E
\]
(EGA IV~8.3.11). For the same reasons, we may suppose that \(T\) is
central above \(U\); it is a torus by 3.6.

Now that \(Q_j\) is known to be ind-constructible, EGA IV~8.3.4
implies that \(Q_j=\operatorname{ens}(S_j)\) for \(j\) sufficiently
large. We may then replace \(S,G,E\) by \(S_j,G_j,E_j\), which
reduces us to the case in which \(S\) is an affine reduced scheme of
finite type over \(\mathbb Z\).

\smallskip
\noindent\textit{c) Reduction to the case in which \(S\) is the
spectrum of a complete reduced noetherian local ring.}\quad
Because the torus with underlying space \(E\) is unique, the usual
technique (EGA IV~8) and Lemma 3.6 allow us to replace \(S\) by the
spectrum of the local ring \(A\) at a point of \(S\). Let
\(\widehat S\) be the spectrum of the completion \(\widehat A\) of
\(A\) for the topology defined by its maximal ideal. Since \(A\) is
the localization of an algebra of finite type over \(\mathbb Z\),
\(\widehat S\) is still reduced (EGA IV~7.6.5).

I claim that it is enough to solve the problem after the base
extension \(\widehat S\to S\). Indeed, let \(\widehat T\) be the
sub-torus of \(G_{\widehat S}\) with underlying space
\(E_{\widehat S}\). Its two inverse images in
\(G_{\widehat S\times_S\widehat S}\) are central sub-tori with the
same underlying space; they therefore coincide
(Exposé \SGARefLink{sga3:ix:theorem:5:1:bis}{IX~5.1 bis}). By fpqc descent, \(\widehat T\) comes from a
torus \(T\) of \(G\) that solves the problem (cf.\ \SGARefLink{sga3:xv:lemma:3:4}{3.4}).

\smallskip
\noindent\textit{d) A descent lemma.}\quad
Recall the following properties of finite morphisms, noted in
TDTE I. Let \(S,S'\) be two preschemes and
\(u:S'\to S\) a finite morphism. Then:
\begin{enumerate}
\item[\normalfont\SGAInlineRefTarget{sga3:xv:lemma:3:6:item:1}(1)] The morphism \(u\) is an epimorphism if and
only if the canonical morphism of sheaves
\[
  \mathcal O_S\longrightarrow u_*(\mathcal O_{S'})
\]
is injective.
\item[\normalfont(2)] The morphism \(u\) is an effective
epimorphism (Exposé \SGARefLink{sga3:iv:definition:1:3}{IV~1.3}) if and only if the canonical diagram
\[
  0\longrightarrow\mathcal O_S
  \longrightarrow u_*(\mathcal O_{S'})
  \rightrightarrows
  u_*(\mathcal O_{S'})\otimes_{\mathcal O_S}u_*(\mathcal O_{S'})
\]
is exact.
\item[\normalfont\SGAInlineRefTarget{sga3:xv:lemma:3:6:item:3}(3)] If in addition \(S\) is noetherian and \(u\)
is an epimorphism, \(u\) is the composite of a finite sequence of
finite effective epimorphisms.
\end{enumerate}

We can now prove the following lemma, whose proof uses a non-flat
descent technique.

\SGAStatement{\SGAInlineRefTarget{sga3:xv:lemma:3-7}Lemma 3.7}{Let \(S\) be a locally noetherian
prescheme, \(G\) an \(S\)-group prescheme of finite type, \(S'\) a
prescheme, and \(u:S'\to S\) a finite morphism. For every
\(S\)-prescheme \(T\), let \(\mathrm M(T)\) be the set of subgroups
of \(G_T\) of multiplicative type \textup{(respectively, the set of
central subgroups of \(G_T\) of multiplicative type)}, so
that \(\mathrm M\) is naturally a contravariant functor on
\(\operatorname{Sch}/S\). If \(u\) is an effective epimorphism
\textup{(respectively, an epimorphism)}, then the diagram of
sets
\[
  \tag{*}\label{sga3:xv:lemma:3-7:equation:star}
  \mathrm M(S)\longrightarrow\mathrm M(S')
  \rightrightarrows\mathrm M(S'\times_S S')
\]
is exact.}
\SGARefTarget{sga3:xv:lemma:3:7}

\SGAProof{Proof}
\textit{i) Reduction to the case in which \(u\) is an effective
epimorphism.}\quad
Here we consider the functor of central subgroups of \(G\) of
multiplicative type. The injectivity of
\(\mathrm M(S)\to\mathrm M(S')\) is local on \(S\); once this
injectivity is granted, the exactness of \((*)\) is also local on
\(S\). We may therefore suppose that \(S\) is affine and
noetherian.

Consider the case in which \(u:S'\to S\) is the composite of two
finite epimorphisms
\[
  S'\xrightarrow{v}S''\xrightarrow{w}S.
\]
I claim that if both diagrams
\[
  \tag{*'}\label{sga3:xv:lemma:3-7:equation:star-prime}
  \mathrm M(S'')\longrightarrow\mathrm M(S')
  \rightrightarrows\mathrm M(S'\times_{S''}S')
\]
and
\[
  \tag{*''}\label{sga3:xv:lemma:3-7:equation:star-prime-prime}
  \mathrm M(S)\longrightarrow\mathrm M(S'')
  \rightrightarrows\mathrm M(S''\times_S S'')
\]
are exact, then so is \((*)\).

Indeed, the injectivity of
\(\mathrm M(S)\to\mathrm M(S')\) is clear. If \(T'\) belongs to
\[
  \operatorname{Ker}\bigl(
    \mathrm M(S')\rightrightarrows
    \mathrm M(S'\times_S S')\bigr),
\]
then, a fortiori, \(T'\) belongs to
\[
  \operatorname{Ker}\bigl(
    \mathrm M(S')\rightrightarrows
    \mathrm M(S'\times_{S''}S')\bigr).
\]
By the exactness of \((*')\), it therefore comes from a unique
element \(T''\) of \(\mathrm M(S'')\). It is enough to show that
\(T''\) belongs to
\[
  \operatorname{Ker}\bigl(
    \mathrm M(S'')\rightrightarrows
    \mathrm M(S''\times_S S'')\bigr),
\]
for then the exactness of \((*'')\) descends \(T''\) to \(S\).
Let \(T''_1,T''_2\) be the two inverse images of \(T''\) in
\(\mathrm M(S''\times_S S'')\). These are central subgroups of
\(G_{S''\times_S S''}\) of multiplicative type, so to prove that
\(T''_1=T''_2\) it is enough to see that they have the same fibers
(Exposé \SGARefLink{sga3:ix:theorem:5:1:bis}{IX~5.1 bis}). Consider the commutative diagram
\[
\begin{tikzcd}[column sep=huge,row sep=large]
S' \arrow[d,"v"'] &
  S'\times_S S' \arrow[l,shift left=.55ex]
    \arrow[l,shift right=.55ex] \arrow[d,"v\times v"]\\
S'' &
  S''\times_S S'' . \arrow[l,shift left=.55ex]
    \arrow[l,shift right=.55ex]
\end{tikzcd}
\]
The morphism \(v\) is a finite epimorphism, hence dominant by
property \textup{\SGARefLink{sga3:xv:lemma:3:6:item:1}{(1)}} above, and closed; it is therefore surjective,
and so is \(v\times v\). Let \(x''\) be a point of
\(S''\times_S S''\), and let \(x'\) be a point of
\(S'\times_S S'\) lying above \(x''\). The commutativity of the
diagram shows that the two inverse images of
\((T''_1)_{x''}\) and \((T''_2)_{x''}\) in \(G_{x'}\) coincide.
Thus
\[
  (T''_1)_{x''}=(T''_2)_{x''}
\]
(EGA IV~2.2.15).

The preceding argument and an immediate induction on the number of
factors in a factorization of \(u:S'\to S\) into effective
epimorphisms, supplied by property \textup{\SGARefLink{sga3:xv:lemma:3:6:item:3}{(3)}} above, show that in
order to prove the exactness of \((*)\) for central subgroups of
multiplicative type it is enough to treat the case in which \(u\)
is an effective epimorphism. Finally, another application of
Exposé \SGARefLink{sga3:ix:theorem:5:1:bis}{IX~5.1 bis} shows that it is enough to prove \SGARefLink{sga3:xv:lemma:3-7}{3.7} when
\(\mathrm M\) is the functor of all subgroups of \(G\) of
multiplicative type and \(u\) is an effective epimorphism.

\smallskip
\noindent\textit{ii) Injectivity of
\(\mathrm M(S)\to\mathrm M(S')\).}\quad
Since \(u:S'\to S\) is an epimorphism, the canonical morphism
\[
  \mathcal O_S\longrightarrow u_*(\mathcal O_{S'})
\]
is injective. Moreover, an \(S\)-group of multiplicative type is
flat over \(S\). The injectivity of
\(\mathrm M(S)\to\mathrm M(S')\) is therefore a consequence of the
following more general lemma.

\SGAStatement{\SGAInlineRefTarget{sga3:xv:lemma:3-8}Lemma 3.8}{Let \(f:X\to S\) and \(g:S'\to S\) be
morphisms of preschemes, let \(\mathcal F\) be a quasi-coherent
\(\mathcal O_X\)-module, and put
\[
  X'=X\times_S S',
  \qquad
  \mathcal F'=\mathcal F\otimes_{\mathcal O_S}\mathcal O_{S'}.
\]
Let \(\mathrm Q(\mathcal F)\) be the set of quotient
\(\mathcal O_X\)-modules \(\mathcal G\) of \(\mathcal F\) that are
quasi-coherent and flat over \(S\), and define
\(\mathrm Q(\mathcal F')\) analogously relative to
\(\mathcal F',X'\), and \(S'\). Suppose that \(g\) is quasi-compact
and that
\[
  \mathcal O_S\longrightarrow g_*(\mathcal O_{S'})
\]
is injective. Then the canonical map
\[
  \begin{aligned}
  \mathrm Q(\mathcal F)&\longrightarrow\mathrm Q(\mathcal F'),\\
  \mathcal G&\longmapsto
    \mathcal G\otimes_{\mathcal O_S}\mathcal O_{S'}
  \end{aligned}
\]
is injective.}
\SGARefTarget{sga3:xv:lemma:3:8}

\SGAProof{Proof}
Indeed, we may suppose that \(S\) is affine, and then that \(X\) is
affine. Since \(S'\) is quasi-compact, it has a finite covering by
affine open subsets. Replacing \(S'\) by the disjoint union of these
open subsets, we may therefore suppose that \(S'\) is affine. We are
then reduced to the following lemma, whose proof is immediate.

\SGAStatement{Lemma 3.9}{Let \(A\to A'\) be an injective
homomorphism of rings, let \(M\) be an \(A\)-module, let \(N=M/P\)
be a quotient \(A\)-module flat over \(A\), and put
\[
  M'=M\otimes_A A',
  \qquad
  N'=N\otimes_A A'=M'/P'.
\]
Then \(P\) is the inverse image of \(P'\) under the canonical
homomorphism \(M\to M'\). Thus \(N\) and \(P\) are determined by
\(N'\) and \(P'\).}
\SGARefTarget{sga3:xv:lemma:3:9}

\smallskip
\noindent\textit{iii) Exactness of \SGARefLink{sga3:xv:lemma:3-7}{3.7} \((*)\) at
\(\mathrm M(S')\).}\quad
Let
\[
  T'\in\operatorname{Ker}\bigl(
    \mathrm M(S')\rightrightarrows
    \mathrm M(S'\times_S S')\bigr).
\]
Suppose that iii) has been proved when \(S\) is the spectrum of a
noetherian local ring. For every point \(s\) of \(S\), there then
exists a subgroup \(T_s\) of multiplicative type in
\(G\times_S\operatorname{Spec}(\mathcal O_{S,s})\), obtained by
descent from
\[
  T'\times_{S'}u^{-1}\!\bigl(\operatorname{Spec}(\mathcal O_{S,s})\bigr).
\]
By Exposé VI\textsubscript{B}, \S10, and \SGARefLink{sga3:xv:lemma:3:6}{3.6}, there is an open
neighborhood \(U\) of \(s\) in \(S\) and a subgroup scheme \(T\) of
\(G\) over \(U\), of multiplicative type, that extends \(T_s\).
Let \(U'\) be the inverse image of \(U\) in \(S'\). We then have two
subschemes of \(G'|_{U'}\), namely \(T'|_{U'}\) and
\(T\times_U U'\), which coincide over
\(u^{-1}(\operatorname{Spec}(\mathcal O_{S,s}))\).

Regarding the latter scheme as the projective limit of the
\(u^{-1}(V)\), where \(V\) runs through the open neighborhoods of
\(s\) in \(U\), EGA IV~8 shows that there is an open subset
\(V\subset U\), containing \(s\), over whose inverse image
\(T\times_U U'\) and \(T'\) coincide. Thus, under the stated
hypotheses, \(T'\) descends locally on \(S\); by the uniqueness
proved in ii), it consequently descends globally. In short, it is
enough to prove iii) when \(S\) is the spectrum of a noetherian
local ring.

Let \(\widehat S\) be the spectrum of the completion of the ring of
\(S\), and put
\[
  S''=S'\times_S S',
  \qquad
  \widehat S'=\widehat S\times_S S',
  \qquad
  \widehat S''=\widehat S\times_S S''
      =\widehat S'\times_S S'.
\]
I claim that it is enough to show that
\[
  \tag{\(\widehat *\)}\label{sga3:xv:lemma:3-9:equation:widehat-star}
  \mathrm M(\widehat S)\longrightarrow
  \mathrm M(\widehat S')
  \rightrightarrows
  \mathrm M(\widehat S'')
\]
is exact at \(\mathrm M(\widehat S')\). This follows from the
commutative diagram
\[
\begin{tikzcd}[column sep=large,row sep=large]
\mathrm M(S) \arrow[r] \arrow[d]
  & \mathrm M(S') \arrow[r,shift left=.55ex]
      \arrow[r,shift right=.55ex] \arrow[d]
  & \mathrm M(S'') \arrow[d]\\
\mathrm M(\widehat S) \arrow[r]
      \arrow[d,shift left=.55ex] \arrow[d,shift right=.55ex]
  & \mathrm M(\widehat S') \arrow[r,shift left=.55ex]
      \arrow[r,shift right=.55ex]
      \arrow[d,shift left=.55ex] \arrow[d,shift right=.55ex]
  & \mathrm M(\widehat S'')\\
\mathrm M(\widehat S\times_S\widehat S)
  \arrow[r,"f"]
  & \mathrm M(\widehat S'\times_{S'}\widehat S').
\end{tikzcd}
\]
Here the second row is exact at \(\mathrm M(\widehat S')\) by
hypothesis, the first two columns are exact by fpqc descent, and
\(f\) is injective by ii), applied to the finite epimorphism
\[
  \widehat S'\times_{S'}\widehat S'
  \longrightarrow
  \widehat S\times_S\widehat S
\]
deduced from \(S'\to S\) by the flat base change
\(\widehat S\times_S\widehat S\to S\). The diagram chase is left
to the reader. The characterization of finite effective
epimorphisms recalled in d)\SGARefLink{sga3:xv:theorem:5:2}{(2)} shows that
\(\widehat S'\to\widehat S\), obtained from \(S'\to S\) by the
flat base change \(\widehat S\to S\), is again a finite effective
epimorphism. We are therefore reduced to the case in which \(S\)
is, in addition, the spectrum of a complete noetherian local ring.

Let \(S_0\) be the reduced closed subscheme of \(S\) whose
underlying space is the closed point. Let \(T'\) be as above and
write \(T''\) for its common image in \(\mathrm M(S'')\). Since
\[
  S'_0=S'\times_S S_0
\]
is faithfully flat over \(S_0\), there is a subgroup \(T_0\) of
\(G_{S_0}\), of multiplicative type, whose inverse image in
\(G_{S'_0}\) is \(T'_{S'_0}\), by fpqc descent. Since \(S\) is
complete noetherian local, there exist an \(S\)-group \(T\) of
multiplicative type and an \(S_0\)-morphism
\[
  u_0:T\times_S S_0\longrightarrow T_0
\]
(Exposé \SGARefLink{sga3:x:corollary:3:3}{X~3.3}). The inverse image of \(u_0\) over \(S'_0\) extends
uniquely to an \(S'\)-isomorphism
\[
  u':T_{S'}\xrightarrow{\sim}T',
\]
again by Exposé \SGARefLink{sga3:x:corollary:3:3}{X~3.3}. Notice here that, because \(S'\) is finite
over the complete local scheme \(S\), it is a finite disjoint sum
of complete local schemes.

The two inverse images of \(u'\) over \(S''\) are morphisms from
\(T_{S''}\) to \(T''\) that coincide over \(S_0\times_S S''\);
hence they coincide everywhere by the same result. Since \(T\) is
flat over \(S\) and \(S'\to S\) is a finite effective
epimorphism, TDTE I, page~8, gives an exact diagram
\[
  \operatorname{Hom}_S(T,G)\longrightarrow
  \operatorname{Hom}_{S'}(T_{S'},G_{S'})
  \rightrightarrows
  \operatorname{Hom}_{S''}(T_{S''},G_{S''}).
\]
Thus \(u'\) comes from a morphism of preschemes \(u:T\to G\).
Likewise, \(T\times_S T\) is flat over \(S\), and therefore
\[
  \operatorname{Hom}_S(T\times_S T,G)\longrightarrow
  \operatorname{Hom}_{S'}(T_{S'}\times_{S'}T_{S'},G_{S'})
\]
is injective. It follows at once that \(u\) is a homomorphism of
groups, since \(u'\) is one.

Moreover, \(u\) is a monomorphism. First observe that
\(\operatorname{Ker}(u)\) is flat over \(S\). To prove this we may
suppose that \(S\) is artinian local (EGA
\(0_{\mathrm{III}}\), 10.2.2), hence that \(G\) is separated
(Exposé VI\textsubscript{A}, \SGARefLink{sga3:VIA:0.3}{0.3}); then
\(\operatorname{Ker}(u)\) is of multiplicative type (Exposé
\SGARefLink{sga3:ix:theorem:6:8}{IX~6.8}), and is therefore flat over \(S\). Since
\[
  \operatorname{Ker}(u)\times_S S'
    =\operatorname{Ker}(u')=0,
\]
\SGARefLink{sga3:xv:lemma:3-8}{Lemma~3.8} gives \(\operatorname{Ker}(u)=0\). Finally, a
monomorphism \(u\) is an immersion (\SGARefLink{sga3:viii:remark:7:13}{Exposé VIII, Remarks~7.13}),
and its image \(u(T)\) is an element of \(\mathrm M(S)\) whose
image in \(\mathrm M(S')\) is \(T'\). This proves \SGARefLink{sga3:xv:lemma:3-7}{Lemma~3.7}.

\smallskip
\noindent\textit{e) End of the proof of I).}\quad
By reduction c), we are reduced to the case in which \(S\) is the
spectrum of a complete reduced noetherian local ring \(A\). Let
\(S'\) be the spectrum of the normalization \(A'\) of \(A\).
Nagata's theorem says that \(A'\) is finite over \(A\) (EGA
\(0_{\mathrm{IV}}\), 23.1.5), and \(S'\to S\) is an epimorphism
because \(A\to A'\) is injective. Suppose that \(G_{S'}\) has a
torus \(T'\) whose underlying space is \(E_{S'}\). The two inverse
images of \(T'\) in \(G_{S'\times_S S'}\) are central sub-tori
with the same underlying space, so they coincide (Exposé IX~5
bis). By \SGARefLink{sga3:xv:lemma:3-7}{Lemma~3.7}, \(T'\) therefore comes from a central
sub-torus \(T\) of \(G_S\), which plainly has \(E\) as its
underlying space. It is thus enough to prove the existence of
\(T'\). This reduces us to the normal case and completes the proof
of I).

\medskip
\noindent\textit{II) Proof that \textup{\SGARefLink{sga3:xv:theorem:3-1:item:iii}{(iii)}} implies
\textup{\SGARefLink{sga3:xv:theorem:3-1:item:ii}{(ii)}} when \(S\) is normal.}\quad
We may restrict ourselves to the case in which \(S\) is integral;
let \(t\) be its generic point. For every integer \(n\) that is a
power of \(q\), let \({}_nG\) be the ``kernel'' in \(G\) of
raising to the \(n\)-th power. Since \(E\) is locally closed in
\(G\), by \SGARefLink{sga3:xv:theorem:3-1:item:iii}{(iii)}(b), the intersection of
\(\operatorname{ens}({}_nG)\) with \(E\) is locally closed in
\(\operatorname{ens}(G)\). Let \(E(n)\) denote the reduced
subscheme of \(G\) whose underlying space is
\(\operatorname{ens}({}_nG)\cap E\).

We show that the structural morphism
\[
  E(n)\longrightarrow S
\]
is separated and universally open. Both properties have a
valuative criterion (EGA II~7.2.3 and EGA IV~14.5.8). Let
\(S'\) therefore be an \(S\)-scheme that is the spectrum of a
complete discrete valuation ring with algebraically closed residue
field. We have proved that \SGARefLink{sga3:xv:theorem:3-1:item:iii}{3.1(iii)} implies \SGARefLink{sga3:xv:theorem:3-1:item:iv}{3.1(iv)}, so there is a
sub-torus \(T'\) of \(G_{S'}\) whose underlying space is
\(E_{S'}\). Now \({}_nT'\) is finite étale over \(S'\), hence is
separated and universally open over \(S'\); the same is therefore
true of \(E(n)\times_S S'\), which has the same underlying space.

The fibers of \(E(n)\), moreover, all have the same number of
geometric points, namely \(n^r\) if \(r\) is the rank of the torus
\(T_t\).\footnote{The French edition prints \(r^n\); the intended
cardinality of the \(n\)-torsion of a rank-\(r\) torus is \(n^r\).}
Finally, the generic fiber \(E(n)_t\), being reduced, is equal to
\({}_nT_t\), and hence is étale over \(\operatorname{Spec}(t)\).
Since \(S\) is normal, SGA I~10.11 now shows that \(E(n)\) is an
étale covering of \(S\).

If \(s\) is a point of \(S\), then \(E(n)_s\) is étale over
\(\operatorname{Spec}\kappa(s)\), hence reduced, and consequently
coincides with the group \({}_nT_s\) of multiplicative type. We
show that \(E(n)\) is a group subscheme of \(G\). Let
\[
  m:E(n)\times_S E(n)\longrightarrow G
\]
be the morphism induced by multiplication in \(G\). The underlying
map \(\operatorname{ens}(m)\) factors through \(E(n)\), so \(m\)
itself factors through the scheme \(E(n)\), because its source
\(E(n)\times_S E(n)\) is étale over \(S\), and hence reduced. It
then follows from Exposé \SGARefLink{sga3:x:corollary:4-8:item:a}{X~4.8(a)} that \(E(n)\) is a subgroup of
multiplicative type. As observed already in \SGARefLink{sga3:xv:remark:3-2:item:a}{3.2(a)}, the family of
subgroups \(E(n)\) is necessarily coherent. This proves that
\SGARefLink{sga3:xv:theorem:3-1:item:iii}{(iii)} implies \SGARefLink{sga3:xv:theorem:3-1:item:ii}{(ii)} when \(S\) is normal.

\medskip
\noindent\textit{III) Proof that \textup{\SGARefLink{sga3:xv:theorem:3-1:item:ii}{(ii)}} implies
\textup{\SGARefLink{sga3:xv:theorem:3-1:item:i}{(i)}}.}\quad
In fact, we shall show that there is a unique sub-torus \(T\) of
\(G\), with underlying space \(E\), such that
\[
  {}_nT=M_n
\]
for every \(n\) that is a power of \(q\). Uniqueness follows
immediately from Exposé \SGARefLink{sga3:ix:corollary:4:8:item:b}{IX~4.8(b)}. In proving existence, in view of
this uniqueness, we may successively suppose that:
\begin{enumerate}
\item[\normalfont a)] \(S\) is noetherian.
\item[\normalfont b)] The \(M_n\) are central subgroups. Indeed, it
is enough to replace \(G\) by
\[
  Z_n=\operatorname{Centr}_G(M_n)
\]
for sufficiently large \(n\), by \SGARefLink{sga3:xv:lemma:2:5}{2.5} and \SGARefLink{sga3:xv:lemma:2-5-bis}{2.5 bis}.
\item[\normalfont c)] \(S\) is the spectrum of a local ring.
Suppose that the problem has been solved after every base change
\(\operatorname{Spec}(\mathcal O_{S,s})\to S\), for \(s\in S\),
and let \(T_s\) be the resulting sub-torus of
\(G\times_S\operatorname{Spec}(\mathcal O_{S,s})\). There is then
an open neighborhood \(U\) of \(s\) and a sub-torus \(T\) of
\(G|_U\) extending \(T_s\), by Exposé VI\textsubscript{B}, \S10,
and \SGARefLink{sga3:xv:lemma:3:6}{3.6}. We have proved that \SGARefLink{sga3:xv:theorem:3-1:item:ii}{3.1(ii)} implies \SGARefLink{sga3:xv:theorem:3-1:item:iv}{3.1(iv)}; since \(S\)
is noetherian, \(E\) is therefore constructible by \SGARefLink{sga3:xv:lemma:3:3}{3.3}. After
shrinking \(U\), we may consequently suppose that
\[
  \operatorname{ens}(T)=E\times_S U
\]
(EGA IV~9.5.2). For every \(n\) that is a power of \(q\), the
groups \({}_nT\) and \(M_n\times_S U\) are then two subgroups of
multiplicative type in \(G|_U\), with the same fibers; because
\(M_n\) is central, they coincide by Exposé \SGARefLink{sga3:ix:corollary:5:3:bis}{IX~5.3 bis}. Thus the
problem has a local solution on \(S\), and uniqueness makes these
local solutions into a global one.
\item[\normalfont d)] \(S\) is the spectrum of a complete
noetherian local ring and, if \(s\) is the closed point, \(T_s\)
is trivial. This follows from EGA \(0_{\mathrm{III}}\), 10.3.1,
and fpqc descent.
\item[\normalfont e)] \(S\) is reduced. Apply \SGARefLink{sga3:xv:corollary:2-2}{2.2}.
\item[\normalfont f)] \(S\) is normal. Apply Nagata's finiteness
theorem (EGA \(0_{\mathrm{IV}}\), 23.1.5) and Lemma~3.7 for
central sub-tori.
\end{enumerate}

After these reductions, \(T_s\) is trivial, and hence
\((M_n)_s\) is trivial. It follows from Exposé \SGARefLink{sga3:x:corollary:3:3}{X~3.3} that \(M_n\)
itself is trivial. If \(t\) is any point of \(S\), all the
subgroups \({}_nT_t\) of \(T_t\), for \(n\) a power of \(q\), are
therefore trivial; Exposé \SGARefLink{sga3:x:proposition:1:4}{X~1.4} then shows at once that \(T_t\)
itself is trivial. It remains only to prove the following lemma.

\SGAStatement{Lemma 3.10}{Under the hypotheses of \SGARefLink{sga3:xv:theorem:3-1:item:ii}{3.1(ii)}, suppose
in addition that \(S\) is noetherian and normal and that, for every
point \(s\) of \(S\), \(T_s\) is a trivial torus. Then there exists
a unique sub-torus \(T\) of \(G\), with underlying space \(E\),
such that
\[
  {}_nT=M_n
\]
for every \(n\) that is a power of \(q\). Moreover, \(T\) is
trivial.}
\SGARefTarget{sga3:xv:lemma:3:10}

\SGAProof{Proof}
Uniqueness follows from the fact that \(T\), being smooth over
\(S\), is reduced. To prove existence, we may suppose that \(S\)
is irreducible, with generic point \(\eta\). Let \(r\) be the rank
of \(T_\eta\), put
\[
  T'=\mathbb G_{m,S}^{\,r},
\]
and choose an isomorphism
\[
  u_\eta:T'_\eta\xrightarrow{\sim}T_\eta.
\]
Write \({}_nu_\eta\) for its restriction to \({}_nT'_\eta\).
Since \({}_nT'\) and \(M_n\) are trivial, \({}_nu_\eta\) has a
unique extension to an \(S\)-isomorphism
\[
  {}_nu:{}_nT'\xrightarrow{\sim}M_n.
\]
I claim that, for every point \(s\) of \(S\), there is a
necessarily unique group isomorphism
\[
  u_s:T'_s\xrightarrow{\sim}T_s
\]
that extends \(({}_nu)_s\) for every \(n\) that is a power of
\(q\).

Indeed, let \(S^1\) be an \(S\)-scheme that is the spectrum of a
complete discrete valuation ring with algebraically closed residue
field, whose generic point \(t_1\) maps to \(\eta\) and whose
closed point \(s_1\) maps to \(s\) (EGA II~7.1.9). The proof that
3.1(ii) implies 3.1(iv) gives a sub-torus \(T^1\) of \(G_{S^1}\)
such that
\[
  (M_n)_{S^1}={}_nT^1
\]
for every \(n\) that is a power of \(q\). Since \(S^1\) is normal,
\(T^1\) is isotrivial by Exposé \SGARefLink{sga3:x:theorem:5-16}{X~5.16}. The classification of
isotrivial tori (Exposé \SGARefLink{sga3:x:corollary:1:2}{X~1.2}) and SGA I, Exposé V~8.2, together
with the fact that the generic fiber of \(T^1\) is trivial, show
that \(T^1\) is trivial. We may therefore extend
\[
  u_\eta\times_\eta t_1
\]
to an \(S^1\)-isomorphism of group schemes
\[
  u^1:T'\times_S S^1\xrightarrow{\sim}T^1.
\]
On \({}_nT'\times_S S^1\), the morphisms \(u^1\) and
\({}_nu\times_S S^1\) agree on the generic fiber, and hence agree
everywhere. Restricting \(u^1\) to the closed fiber gives the
required \(u_s\) after the field extension
\(\kappa(s)\to\kappa(s_1)\); it descends to \(\kappa(s)\) by
Exposé \SGARefLink{sga3:ix:corollary:4:8:item:a}{IX~4.8(a)} and fpqc descent.

The scheme \(T'\) is smooth over the normal scheme \(S\), hence is
normal. The extension criterion for rational maps in EGA
IV\textsubscript{4}, 20.4.6, now shows that there is an
\(S\)-morphism
\[
  u:T'\longrightarrow G
\]
whose restriction to each fiber \(T'_s\) is the composite
\[
  T'_s\xrightarrow{u_s}T_s\longrightarrow G_s.
\]
Indeed, for every base change \(S^1\to S\) of the kind used above,
the preceding construction gives precisely the required extension
over \(S^1\).

It remains to show that \(u\) is a homomorphism. Let \(m_{T'}\)
and \(m_G\) denote multiplication in \(T'\) and \(G\),
respectively. The coincidence subscheme of
\[
  m_G\circ(u\times_Su)
  \quad\text{and}\quad
  u\circ m_{T'}
\]
contains every fiber, since every \(u_s\) is a homomorphism. It
therefore has the same underlying space as \(T'\times_S T'\), and
so equals \(T'\times_S T'\), because the latter is smooth over
\(S\) and hence reduced. Finally, \(u\) is a monomorphism because
it is one on every fiber; it is therefore an immersion (Exposé
\SGARefLink{sga3:viii:proposition:7:9}{VIII~7.9}). The image of \(T'\) under \(u\) is a sub-torus of \(G\)
whose underlying set is \(E\).

This completes the proof of \SGARefLink{sga3:xv:theorem:3-1}{Theorem~3.1}.
